(1) 均匀球体内半径为 r 的包围质量 \(M(r)=Mr^3/R^3\)。球对称壳层定理给
\[
\boldsymbol g(r)=-\frac{GM(r)}{r^2}\hat{\boldsymbol r}
=-\frac{GM}{R^3}\boldsymbol r,\quad 0\le r\le R.
\]
以无穷远引力势能为零、在球面连续，积分得
\[
U(r)=-\frac{GMm}{2R}\left(3-\frac{r^2}{R^2}\right).
\]
球内力正比位移，是三维各向同性谐振力场。

(2) 斜弦到球心的距离 d 固定，沿隧道中心的坐标 s 对应 \(r^2=d^2+s^2\)。将引力投影到隧道方向，得
\[
m\ddot s=-\frac{GMm}{R^3}s,\quad
\omega=\sqrt{GM/R^3}.
\]
弦端坐标为 \(s=\pm a\)，其中 \(a=\sqrt{R^2-d^2}\)。从一端静止释放可选 \(s(0)=a,\dot s(0)=0\)，故 \(s(t)=a\cos\omega t\)。到达另一端 \(s=-a\) 的首次时刻为
\[
t_{\rm traverse}=\pi/\omega=\pi\sqrt{R^3/(GM)}.
\]

(3) 机械能守恒，沿 s 的有效势差为 \(\Delta U=m\omega^2(s^2-a^2)/2\)，得到
\[
\dot s^2=\omega^2(a^2-s^2),\quad
|\dot s(s)|=\omega\sqrt{R^2-d^2-s^2}.
\]
最大速度出现在隧道中点 s=0，为 \(v_{\max}=\sqrt{GM/R^3}\sqrt{R^2-d^2}\)。随着 d 接近 R，隧道缩短，最大速度趋零。

(4) 对 d=0 的直径隧道，\(a=R\)、\(v_{\max}=\sqrt{GM/R}\)，穿越时间仍为 \(\pi /\omega \)。对 \(d\to R^-\)，长度 2a 与最大速度都趋零，但每个非退化隧道的首次另一端时间仍恰为 \(\pi /\omega \)；严格 d=R 时隧道退化为一个点，谈穿越时间没有意义。时间的普适性来自均匀天体内引力恰是线性回复力，振幅变化不改变理想谐振周期。
